\section{Stopping the expansion at its first excluded prefix}
\label{p1:stop:section}

The leading correction comes from retaining one dependent pair.
An expansion of the entire counting law would also retain every
possible interaction after that pair.  Taking absolute values at
that stage can be expensive, even when the interactions that first
leave the retained class have small total probability.  We therefore
stop each term at its first excluded prefix and leave all subsequent
indicators inside a probability kernel.

The first identity is finite and does not require independence.
Its exact-marginal application will reduce the approximation problem to
families whose effective nullity is at most $2L$.
The first-excluded-prefix identity is a specialization of
Zeilberger's abstract lace expansion
\cite[formula~(*)]{zeilberger1997}. The broader inclusion--exclusion
context includes the abstract-tube identities
\cite{naimanwynn1997,dohmen1999}. We give the prefix proof and its
convolution form; the model-specific estimates concern the retained
arithmetic class and its full boundary.

\subsection{An identity for an arbitrary hereditary class}

Let $G$ be a finite totally ordered set.  A family
$\mathscr D\subseteq 2^G$ is hereditary if $I\in\mathscr D$ and
$J\subseteq I$ imply $J\in\mathscr D$.  Assume throughout that
$\varnothing\in\mathscr D$.  Its \emph{ordered boundary} is
\begin{equation}
 \partial\mathscr D
 =\{I\subseteq G:I\notin\mathscr D,\ 
           I\setminus\{\max I\}\in\mathscr D\}.
 \label{p1:stop:boundary}
\end{equation}
This definition records the first excluded prefix in the prescribed
order.  It is not the family of all inclusion-minimal excluded sets:
an element of $\partial\mathscr D$ can contain a smaller excluded
subset which is not a prefix.

\begin{lemma}[Expansion at the first excluded prefix]
\label{p1:stop:polynomial}
Let $(u_x)_{x\in G}$ belong to a commutative unital algebra.
Then, as a polynomial identity,
\begin{equation}
 \prod_{x\in G}(1+t u_x)
 =\sum_{I\in\mathscr D}t^{|I|}\prod_{x\in I}u_x
  +\sum_{I\in\partial\mathscr D}
       t^{|I|}\prod_{x\in I}u_x
       \prod_{\substack{y\in G\\y>\max I}}(1+t u_y).
 \label{p1:stop:polynomial-identity}
\end{equation}
\end{lemma}

\begin{proof}
Expand the product on the left and index its monomials by
$S\subseteq G$.  If $S\notin\mathscr D$, read the elements of $S$
in their prescribed order.  There is a unique first prefix
$I\notin\mathscr D$.  Its preceding prefix belongs to $\mathscr D$,
so $I\in\partial\mathscr D$, and every element of $S\setminus I$
is larger than $\max I$.

Conversely, take $I\in\partial\mathscr D$ and
$T\subseteq\{y\in G:y>\max I\}$.  Every proper prefix of $I$
belongs to $\mathscr D$, by heredity.  The set $I\cup T$ is
excluded, since it contains $I$, and its first excluded prefix is
$I$.  Thus the pairs $(I,T)$ partition the monomials indexed by
$S\notin\mathscr D$.  Summing these monomials proves the identity.
\end{proof}

To identify the specialization explicitly, define \(\ell(S)=S\)
for \(S\in\mathscr D\), and otherwise let \(\ell(S)\) be its
first excluded prefix. This is a lace map: \(\ell(S)\subseteq S\);
\(\ell(T)=\ell(S)\) whenever \(\ell(S)\subseteq T\subseteq S\);
and equal images of two sets remain unchanged under their union.
Its fixed points are \(\mathscr D\cup\partial\mathscr D\).
The compatible-extension set is empty for \(I\in\mathscr D\),
and equals \(\{y\in G:y>\max I\}\) for
\(I\in\partial\mathscr D\). Zeilberger's formula~(*) with
\(Y_x=t u_x\) is therefore exactly
\eqref{p1:stop:polynomial-identity}.

For indicators $(J_x)_{x\in G}$, write
\[
 Z=\sum_{x\in G}J_x,\qquad
 q(I)=\E\prod_{x\in I}J_x,\qquad q(\varnothing)=1.
\]
Let $\Delta=\delta_1-\delta_0$, and use the convention
$\Delta^{*0}=\delta_0$.  The measure retained by $\mathscr D$ is
\begin{equation}
 T_{\mathscr D}
       =\sum_{I\in\mathscr D}q(I)\Delta^{*|I|}.
 \label{p1:stop:retained}
\end{equation}
It has total mass one, although its individual masses need not
be nonnegative.

\begin{proposition}[A positive remainder kernel]
\label{p1:stop:measure}
For arbitrary dependence among the indicators,
\begin{align}
 \Law Z-T_{\mathscr D}
   &=\sum_{I\in\partial\mathscr D}\Delta^{*|I|}*\nu_I,
       \label{p1:stop:measure-identity}\\
 \nu_I
   &=\E\left[
        \left(\prod_{x\in I}J_x\right)
        \delta_{\sum_{y>\max I}J_y}
        \right].
       \notag
\end{align}
The measure $\nu_I$ is positive and has mass $q(I)$.  Consequently,
\begin{equation}
 \|\Law Z-T_{\mathscr D}\|_1
       \le \sum_{I\in\partial\mathscr D}2^{|I|}q(I).
 \label{p1:stop:norm}
\end{equation}
\end{proposition}

\begin{proof}
Apply Lemma~\ref{p1:stop:polynomial} in the convolution algebra
of finite signed measures, with unit $\delta_0$ and $t=\Delta$.
For an indicator,
\[
 \delta_0+J_x\Delta=\delta_{J_x},
 \qquad
 \mathop{*}_{y>\max I}(\delta_0+J_y\Delta)
       =\delta_{\sum_{y>\max I}J_y}.
\]
The expectation of the complete product is $\Law Z$.
Taking expectations in \eqref{p1:stop:polynomial-identity}
therefore gives \eqref{p1:stop:measure-identity}.
Positivity and the mass formula for $\nu_I$ follow directly
from its definition.  Finally,
$\|\Delta^{*|I|}\|_1=2^{|I|}$ and the convolution norm is
submultiplicative, which proves \eqref{p1:stop:norm}.
\end{proof}

If $q(I)>0$, the normalized measure $\nu_I/q(I)$ is the law
of the suffix count conditioned on the activity of every member
of $I$.  This conditional law can differ substantially from its
unconditional version.  The proof uses only its total mass,
so no independence between the prefix and suffix is required.
In particular, it incurs no factor exponential in the number
of suffix indicators.
The use of the ordinary $\ell^1$ norm matters: an exponential
weight would also charge the suffix count.

\begin{example}[A suffix that remains fully dependent]
\label{p1:stop:example}
Take $G=\{1,2,3\}$ in its natural order, and retain only the empty
set and the singletons.  Then
\[
 \partial\mathscr D=\{\{1,2\},\{1,3\},\{2,3\}\},
 \qquad
 T_{\mathscr D}=\delta_0+(q_1+q_2+q_3)\Delta,
\]
and the exact remainder is
\[
 \Law Z-T_{\mathscr D}
 =\Delta^{*2}*
   \left(\E[J_1J_2\delta_{J_3}]
               +(q_{13}+q_{23})\delta_0\right).
\]
For $J_1=J_2=J_3=H$, where $\P(H=1)=\theta$, the conditional
suffix following $\{1,2\}$ is identically one.  The formula becomes
\[
 \Law(3H)-(\delta_0+3\theta\Delta)
    =\theta\Delta^{*2}*(2\delta_0+\delta_1).
\]
The remainder has norm $6\theta$, while
\eqref{p1:stop:norm} gives the valid bound $12\theta$.
The mechanism is the positive kernel of mass $3\theta$, not
a comparison of the conditional suffix with an independent one.
This finite example explains the stopping identity; the remaining
task in our model is to prove that the appropriate boundary has
small weighted mass.
\end{example}
